\subsection{An exact comparison with the consecutive repeated product}
\label{sec:kt-cubic-shuffle}

Retain all the data and local conditions of
KT-ORDINARY-CUBIC-MASSEY-CONDITION.  All cochains use
\(df(g,h)=f(g)+g f(h)-f(gh)\), and the cup convention is
\((x\smile y)(g,h)=x(g)\,g y(h)\).
Pointwise multiplication of one-cochains is denoted by
juxtaposition.  In particular it is not the cup product.

\begin{proposition}[KT-ORDINARY-CUBIC-SHUFFLE]
Define
\[
 h=-\frac12u^2,\qquad a=-w-\frac12uv,\qquad
 b=w-\frac12uv .                                      \tag{CS.1}
\]
Then \(dh=u\smile u\), \(da=u\smile v\), and
\(db=v\smile u\).  The cocycle
\[
 \Phi=u\smile a+h\smile v\in Z^2(G_{\mathbf Q,S},k(\psi))
                                                               \tag{CS.2}
\]
is a specified consecutive defining value for
\(\langle u,u,v\rangle\).  With
\(\Theta=u\smile b+a\smile u\) as in (CM.20), the exact
cochain identities are
\[
 \Theta+2\Phi=-d(ua),\qquad
 \Omega_+=3\Phi-dR,\qquad
 R=\frac32uw+\frac12u^2v.                              \tag{CS.3}
\]

If \((c,(z_v)_v)\), with \(c=-\Omega_+\), represents
the ordinary obstruction in (CM.14), put
\[
 z^{\mathrm{rep}}_v=\frac{R|_{G_{\mathbf Q_v}}-z_v}{3}.
                                                               \tag{CS.4}
\]
Then \(d z^{\mathrm{rep}}_v=\Phi|_{G_{\mathbf Q_v}}\),
and the relative class
\[
 \mathfrak m_{\mathrm{rep}}(u,u,v)
       =[(\Phi,(z^{\mathrm{rep}}_v)_v)]\in\mathcal O_\psi
                                                               \tag{CS.5}
\]
is well defined under the same permitted global and local
equivalences as the ordinary cubic obstruction.  Exactly
\[
             \mathfrak o_3(u,v)=-3\mathfrak m_{\mathrm{rep}}(u,u,v).
                                                               \tag{CS.6}
\]
Thus changing the order of the repeated entries introduces
no additional gap for this specified relative obstruction.
\end{proposition}

\begin{proof}
The cocycle identity for \(u\) gives
\(d(u^2)=-2u\smile u\).  The other two defining equations
are (CM.19).  The differential and cup rule show that the
two terms in \(d\Phi\) cancel, so \(\Phi\) is a cocycle.

For a direct sign check use the six variables of (CM.9).
They give
\[
\begin{split}
 \Phi(g,h)&=-Ac-\frac12Aab-\frac12A^2b,\\
 \Theta(g,h)&=Ac-Ca-\frac12Aab-\frac12ABa,\\
 d(uw)(g,h)&=-Ac-Ca-\frac12(A+a)(Ab-Ba),\\
 d(u^2v)(g,h)&=-A^2b-2ABa-2Aab-Ba^2.
\end{split}                                                   \tag{CS.7}
\]
Substitution proves the first identity in (CS.3), since
\(-ua=uw+u^2v/2\).  Substituting that identity in (CM.21)
gives the second identity, with precisely the displayed \(R\).
All denominators are units because \(p>3\).

For the relative assertion, (CS.3) gives
\(c=-3\Phi+dR\), and (CS.4) is therefore a local primitive
of \(\Phi\).  At the level of pairs one has the literal equality
\[
 (c,(z_v)_v)=-3(\Phi,(z^{\mathrm{rep}}_v)_v)
                              +(dR,(R|_v)_v).                 \tag{CS.8}
\]
The last pair is zero in (CM.14).  Since the allowed local
spaces \(L_v\) are vector subspaces and \(3\) is invertible,
any change by an allowed local cocycle transports to an
allowed local cocycle.  A change by a global primitive
transports to the corresponding global primitive.
The already proved independence of \(\mathfrak o_3\),
together with (CS.8), proves the asserted independence
of (CS.5).  This transports the actual ordinary local
choices; it does not replace them by arbitrary local
primitives for an unrestricted Massey product.
\end{proof}

\paragraph{The exact relation to Sharifi's conventions.}
Here is a further algebraic comparison, separate from the
proper norm-system requirement in Sharifi's theorem.
Restrict to \(K=\mathbf Q(\sqrt D,\mu_p)\), so the
coefficient characters become trivial, and first take
\(k=\mathbf F_p\).  A fixed primitive \(p\)-th root
identifies the coefficients with \(\mu_p\).
In the upper-unitriangular four-by-four group modulo its
central \((1,4)\) entry, the functions
\[
 \kappa_{12}=\kappa_{23}=u,\quad\kappa_{34}=v,\quad
 \kappa_{13}=-h=\frac12u^2,\quad\kappa_{24}=-a
                                                               \tag{CS.9}
\]
give a homomorphism.  Indeed Sharifi's defining equation
(his condition (c), printed p.~18) is
\(d\kappa_{ij}=-\sum_l\kappa_{il}\smile\kappa_{lj}\),
which is exactly (CS.1) with these signs.
Conjugation by \(I-E_{12}/2\) changes only the noncentral
entry \(\kappa_{13}\), replacing it by
\(\binom{u}{2}=(u^2-u)/2\).  Its values therefore lie
in the subgroup \(G_1^{(2)}/Z_1^{(2)}\) generated in
Sharifi's notation by the images of
\(X=I+E_{12}+E_{23}\) and \(Y=I+E_{34}\).
There is no remaining discrepancy in the repeated-entry
block or its binomial normalization.

With the convention (4.1) of that paper, the cocycle
of this last defining system is
\[
 \kappa_{12}\smile\kappa_{24}
       +\kappa_{13}\smile\kappa_{34}
       =-\Phi-\frac12d a .                              \tag{CS.10}
\]
Thus its raw class is \(-[\Phi]\), and the raw ordinary
cubic obstruction is three times this class.  The same
global-coboundary transport as (CS.8) adjusts the local
primitives.  For larger finite coefficient fields the
cochain identities remain valid over that field; a
literal application of a theorem with \(\mu_p\)
coefficients must still identify the requisite scalar
components rather than treating a general \(k\)-valued
character as one cyclic degree-\(p\) extension.

\paragraph{What this comparison does and does not supply.}
Sharifi defines a \emph{proper} system more restrictively:
for the cyclic Kummer extension \(L/K\) attached to the
first entry, it must come from an element \(y\in L^\times\)
with prescribed norm and with \(D^{(1)}y\) admitting a
\(p\)-th root in the fixed extension unramified outside
the chosen set \(S\).  This is the definition following
Proposition~4.2, printed p.~21.  The subgroup and sign
identifications (CS.9)--(CS.10) do not by themselves
prove that our particular ordinary-normalized cochain
is obtained from such a \(y\).  The chosen local
primitives would also need to be matched.

If that matching is established, Theorems~4.3 and 5.4
give, in the case \(n=m=1\), \(k=2<p-1\), a norm-ideal
class formula for the raw value, modulo the stated
lower-product indeterminacy.  The ordinary relative
obstruction retains the local correction in (CM.15),
which is not determined by a raw class modulo that
indeterminacy.  No norm-ideal class has been evaluated
here, and neither its nonvanishing nor injectivity in
the original strict class has been proved.  Consequently
(CS.6) is an exact equivalence of obstruction expressions,
not a new arithmetic exclusion or a source-to-jet theorem.

The primary passages inspected for this comparison are
Sharifi, \emph{Massey products and ideal class groups},
author version dated 26 July 2006: Section~4, printed
pp.~17--22, especially condition (c), (4.1), the two
generators, Proposition~4.2 and the proper-system
definition; and Theorem~5.4 with its complete proof,
printed pp.~26--27.  The published article is
\emph{J.~reine angew.~Math.}~603 (2007), 1--33.
